\documentclass[aspectratio=169]{beamer}

\usetheme{Berkeley}
\usecolortheme{default}

\usepackage[T1]{fontenc}
\usepackage[utf8]{inputenc}
\usepackage{lmodern}
\usepackage{amsmath}
\usepackage{booktabs}

\title{MA2008B Exam Instructions}
\subtitle{How to organize your work and how the exam will be graded}
\author{MA2008B}
\date{}

\setbeamertemplate{navigation symbols}{}
\setbeamertemplate{itemize items}[circle]
\setbeamertemplate{enumerate items}[default]

\begin{document}

\begin{frame}
  \titlepage
\end{frame}

\begin{frame}{Exam Format}
  \begin{itemize}
    \item The exam contains \textbf{five problems}.
    \item You must select \textbf{exactly four problems} to be graded.
    \item The maximum time allowed is \textbf{90 minutes}.
    \item Only the problems you select will be graded.
  \end{itemize}

  \vspace{0.6cm}

  \begin{block}{Before you start}
    Mark the four selected problem boxes clearly on the first page of the exam.
  \end{block}
\end{frame}

\begin{frame}{Problem Selection Checklist}
  On the exam cover page, mark exactly four boxes:

  \vspace{0.4cm}

  \begin{center}
    \begin{tabular}{ccccc}
      $\square$ Q1 & $\square$ Q2 & $\square$ Q3 & $\square$ Q4 & $\square$ Q5
    \end{tabular}
  \end{center}

  \vspace{0.6cm}

  \begin{alertblock}{Important}
    If you solve more than four problems, only the four selected problems are guaranteed to be graded.
  \end{alertblock}
\end{frame}

\begin{frame}{How to Organize Your Solutions}
  \begin{itemize}
    \item Start each selected problem on a \textbf{new page}.
    \item Do not write two different problem solutions on the same page.
    \item Write the selected problems in \textbf{increasing question order}.
    \item If a solution uses several pages, number them clearly:
      \[
        \text{Page 1 of 3}, \quad \text{Page 2 of 3}, \quad \text{Page 3 of 3}.
      \]
  \end{itemize}
\end{frame}

\begin{frame}{Problem Names and Final Answers}
  At the top center of the first page of each selected solution, write the problem name exactly as it appears in the exam.

  \vspace{0.5cm}

  \begin{exampleblock}{Example}
    \centering
    Question 3 -- 5.3 Fundamental Matrix Computation
  \end{exampleblock}

  \vspace{0.5cm}

  For every selected problem:
  \begin{itemize}
    \item include the main computations or reasoning;
    \item clearly state the final answer, classification, or conclusion.
  \end{itemize}
\end{frame}

\begin{frame}{Why Reasoning Matters}
  Unsupported final answers may not be counted as correct.

  \vspace{0.4cm}

  A strong solution usually includes:
  \begin{itemize}
    \item the equations or criteria you are using;
    \item the key algebraic or computational steps;
    \item a clear interpretation of the result;
    \item the final answer written in a way that matches the question.
  \end{itemize}
\end{frame}

\begin{frame}{Grading Rule for Four Selected Problems}
  Correct selected problems are weighted marginally:

  \vspace{0.3cm}

  \begin{center}
    \begin{tabular}{cc}
      \toprule
      Correct selected problems & Final score \\
      \midrule
      0 & 0/10 \\
      1 & 4/10 \\
      2 & 7/10 \\
      3 & 9/10 \\
      4 & 10/10 \\
      \bottomrule
    \end{tabular}
  \end{center}

  \vspace{0.3cm}

  The first correct problem is worth 4 points, the second 3 points, the third 2 points, and the fourth 1 point.
\end{frame}

\begin{frame}{Partial Credit}
  Partial credit may be awarded for partially correct selected problems.

  \vspace{0.4cm}

  Fractional correctness is assigned to maximize your score using the same marginal weights:
  \[
    4,\quad 3,\quad 2,\quad 1.
  \]

  \vspace{0.4cm}

  \begin{block}{Awarded points}
    \[
      \text{awarded points} =
      \text{progress percentage} \times \text{maximum points}.
    \]
  \end{block}
\end{frame}

\begin{frame}{Worked Grading Example}
  Suppose your four selected problems are graded as:
  \[
    1,\quad \frac{2}{3},\quad \frac{1}{3},\quad 0.
  \]

  \vspace{0.3cm}

  Then your score is computed as:
  \[
    4
    + 3\left(\frac{2}{3}\right)
    + 2\left(\frac{1}{3}\right)
    + 1(0)
    = \frac{20}{3}
    \approx 6.67.
  \]

  \vspace{0.3cm}

  So the final score is approximately \textbf{6.67/10}.
\end{frame}

\begin{frame}{Before Submitting}
  Check that:

  \begin{itemize}
    \item you selected exactly four problems;
    \item each selected problem starts on a new page;
    \item each selected problem has the correct problem name at the top;
    \item multi-page solutions are numbered clearly;
    \item your reasoning and computations are visible;
    \item every selected problem has a final answer or classification.
  \end{itemize}
\end{frame}

\begin{frame}{Main Takeaway}
  \begin{center}
    \Large
    Select exactly four problems, organize them clearly, show your reasoning, and state your final answers.
  \end{center}

  \vspace{0.8cm}

  \begin{center}
    \normalsize
    Clear structure makes it easier to grade your work correctly.
  \end{center}
\end{frame}

\end{document}
